\documentclass[a4paper]{article}
\usepackage{tikz}
\usetikzlibrary{shapes.geometric, shadows,calc, layers}

\begin{document}

\section*{TikZ-Beispiel: Platzierung von Nodes mit Referenzpunkten}

\begin{tikzpicture}[remember picture,overlay]
  \draw[help lines] (current page.north west) grid (current page.south
  east);
  \fill[red] (current page.north east) circle (5pt);
  \fill[blue] ([xshift=-50mm]current page.north east) circle (5pt);
  % Ihre Zeichnung hier
\end{tikzpicture}

\begin{tikzpicture}%[remember picture,overlay]
  \coordinate (A) at (0,0);
  \coordinate (B) at ($(A) + (2cm,1cm)$);
  
  \draw[fill=blue!20] (A) circle (3pt) node[below left] {A};
  \draw[fill=red!20] (B) circle (3pt) node[above right] {B};
  \draw[->] (A) -- (B) node[midway, above] {Bewegung};
\end{tikzpicture}

\begin{tikzpicture}

%  \coordinate (X) at (3,5);
%  \fill[red] (X) circle (2pt); % Roter Punkt an Referenz
  
% 1. Ein einfacher Node mit Text und einem Rechteckrahmen
\node[draw] at (0, 0) (node1) {Ein einfacher Textblock};
\fill[red] (node1) circle (2pt); % Roter Punkt an Referenz

%\node[draw] at ([xshift=-25mm,yshift=-70mm]current page.north east) (node12) {Ein einfacher Textblock};
%\fill[red] (current page.north east) circle (2pt); % Roter Punkt an Referenz

% 2. Node mit abgerundeten Ecken und Hintergrundfarbe
\node[draw, rounded corners, fill=blue!20] at (6, 0) (node2) {Abgerundete Ecken und Farbe};
\fill[red] (node2) circle (2pt); % Roter Punkt an Referenz

% 3. Node mit dickem Rahmen und Schatten
\node[draw, very thick, drop shadow] at (0, -2) (node3) {Dicker Rahmen mit Schatten};
\fill[red] (node3) circle (2pt); % Roter Punkt an Referenz

% 4. Node mit geometrischer Form (Ellipse)
\node[draw, ellipse, fill=green!20] at (6, -2) (node4) {Elliptischer Rahmen};
\fill[red] (node4) circle (2pt); % Roter Punkt an Referenz

% 5. Node mit fester Breite und zentriertem Text
\node[draw, text width=4cm, align=center] at (0, -4) (node5) {Feste Breite und zentrierter Text};
\fill[red] (node5) circle (2pt); % Roter Punkt an Referenz

% 6. Node mit benutzerdefiniertem Stil
\tikzset{customstyle/.style={draw, fill=yellow!50, text width=3cm, align=center, rounded corners}}
\node[customstyle] at (6, -4) (node6) {Node mit benutzerdefiniertem Stil};
\fill[red] (node6) circle (2pt); % Roter Punkt an Referenz

% 7. Zwei Nodes mit Verbindungslinie
\node[draw, circle] at (-3, -6) (node7) {Node A};
\fill[red] (node7) circle (2pt); % Roter Punkt an Referenz
\node[draw, circle] at (3, -6) (node8) {Node B};
\fill[red] (node8) circle (2pt); % Roter Punkt an Referenz
\draw[->] (node7) -- (node8);

% \node[
%   draw=red,        % Zeichnet einen roten Rahmen
%   thick,           % Setzt die Linienstärke auf 'dick'
%   inner sep=5pt,   % Abstand zwischen Inhalt und Rahmen
%   anchor=north east,
%   align=right
% ] at ([xshift=-25mm,yshift=-70mm]current page.north east) {%
%   \begin{tabular}{r}
%     \textbf{Praxis Mustermann} \\
%     Musterstraße 123 \\
%     12345 Musterstadt \\
%     Telefon: 01234 / 567890 \\
%     Email: praxis@example.com \\
%   \end{tabular}
% };


\end{tikzpicture}



% \tikzset{my doodle/.pic=%
%     {%
%         \coordinate (-A) at (0,0);
%         \coordinate (-B) at (1cm, 0);
%         \coordinate (-C) at (1cm, -1cm);
%         \coordinate (-D) at (0, -1cm);
%         \draw[line width=1.5mm,red!40!orange!30!yellow] (-A) -- (-B) -- (-C) -- (-D) -- cycle; % gold rectangle
%     % blue line
        
%         \coordinate (-D1) at (-0.055cm,-1.07cm);
%         \path let \p1=(-D1) in coordinate (-D2) at (-0.25cm,\y1+0.1cm) ; 
%         \path let \p2=(-D2) in coordinate (-D3) at (-0.1cm,\y2-0.05cm); 
        
%         \draw[blue] (-D1) -- (-D2) -- (-D3);
%     }%
% }

\end{document}
